problem:  
Let \((N^{2n}, J, h)\) be a complete Kähler manifold with **strongly seminegative holomorphic bisectional curvature** in the sense of Siu.  
For each \(q \in N\), define the curvature nullity subspace
\[
\mathcal{N}_q := \{ X \in T^{1,0}_q N \mid R^N(X,\overline{X},Y,\overline{Y}) = 0 \text{ for all } Y \in T^{1,0}_q N \},
\]
and assume that these assemble into a smooth, \(J\)-invariant, holomorphic distribution \(\mathcal{N} \subset T^{1,0}N\), locally integrable to totally geodesic complex leaves.

Let \(f_0:(M^m,g)\to (N,h)\) be a smooth **energy-minimizing harmonic map** from a compact Riemannian manifold \(M\).  
Let \(V \in \Gamma(f_0^*TN)\) be a **Jacobi field** along \(f_0\); i.e. \(V\) satisfies the linearized harmonic map equation
\[
J_{f_0}(V) := \Delta V + \operatorname{trace}_g R^N(V,df_0(\cdot))\,df_0(\cdot) = 0.
\]

**Conjectural problem:**  
Under the curvature assumption above, must every Jacobi field \(V\) associated to an *energy-minimizing* harmonic map \(f_0\) satisfy
\[
V(p) \in \mathcal{N}_{f_0(p)} \quad \text{for all } p \in M?
\]
Equivalently, does the nullity distribution \(\mathcal{N}\) govern all infinitesimal harmonic deformations of energy-minimizing maps—so that the **kernel of the Jacobi operator** coincides with the sections of \(f_0^*\mathcal{N}\)?  

Furthermore, if such a vanishing holds, does it imply that the **local moduli space** of energy-minimizing harmonic maps near \(f_0\) is modeled on the space of leafwise holomorphic sections of \(f_0^*\mathcal{N}\)?

---

Why is it a "good" problem:  
1. **Precise analytic structure:**  
   This version articulates the problem in terms of the Jacobi operator, the canonical analytic linearization governing harmonic map deformations. It asks whether the Bochner term associated with strongly seminegative curvature forces the *kernel* of this elliptic operator—i.e., infinitesimal deformations—to be tangent to the curvature nullity distribution.  

2. **Natural extension of vanishing theorems:**  
   Classical rigidity results (Siu, Sampson) show that under strong seminegativity, certain components of the second fundamental form vanish, leading to pluriharmonic or holomorphic rigidity. The present conjecture extends this vanishing principle to **the deformation level**, probing whether curvature degeneracy completely confines infinitesimal harmonic moduli.  

3. **Bridge between fields:**  
   The question connects *geometric analysis* (Jacobi fields, Bochner identities) with *complex differential geometry* (nullity distributions, holomorphic foliations) and *deformation theory* (kernel–moduli correspondence). Any progress would deepen the analytic foundation of vanishing theorems and illuminate the geometry of harmonic map moduli spaces.  

4. **Simplicity with depth:**  
   The formulation is concise—an explicit condition on Jacobi fields—but conceptually profound: it asks whether curvature degeneracy not only restricts maps but also freezes their infinitesimal deformations. A resolution would clarify the geometric meaning of strong seminegativity in the analytic and moduli-theoretic context, potentially creating a new aspect of harmonic rigidity theory.  

Hence, this is a *good* research-level problem: simple to state yet rich in analytic, geometric, and structural consequences, extending vanishing phenomena from individual harmonic maps to their entire infinitesimal deformation space.